\PassOptionsToPackage{unicode}{hyperref}
\PassOptionsToPackage{hyphens}{url}
\documentclass[11pt,a4paper]{amsart}
\providecommand{\tightlist}{\setlength{\itemsep}{0pt}\setlength{\parskip}{0pt}}
\usepackage{amsthm}
\usepackage{amsfonts}
\usepackage{amsmath,amssymb}
\usepackage{mathrsfs}
\usepackage{mathtools}
\usepackage{stmaryrd}
\usepackage{bbm}
\usepackage{enumitem}
\usepackage{slashed}
\usepackage{tikz-cd}
\usepackage{tikz}
\providecommand{\Tr}{\mathrm{Tr}}
\providecommand{\Sym}{\mathrm{Sym}}
\providecommand{\Spec}{\mathrm{Spec}}
\providecommand{\Hom}{\mathrm{Hom}}
\providecommand{\End}{\mathrm{End}}
\providecommand{\Gal}{\mathrm{Gal}}
\providecommand{\Pic}{\mathrm{Pic}}
\providecommand{\NS}{\mathrm{NS}}
\providecommand{\rank}{\mathrm{rank}}
\providecommand{\disc}{\mathrm{disc}}
\providecommand{\OK}{\mathcal{O}_K}
\providecommand{\K}{\mathbb{K}}
\providecommand{\Z}{\mathbb{Z}}
\providecommand{\Q}{\mathbb{Q}}
\providecommand{\C}{\mathbb{C}}
\providecommand{\fp}{\mathfrak{p}}
\providecommand{\frob}{\mathrm{Frob}}

\usepackage[T1]{fontenc}
\usepackage[utf8]{inputenc}
\usepackage{lmodern}
\usepackage{xcolor}
\usepackage[margin=1in]{geometry}
\usepackage{longtable,booktabs,array}
\usepackage{microtype}
\usepackage{hyperref}
\hypersetup{hidelinks}

\newtheorem{theorem}{Theorem}[section]
\newtheorem{lemma}[theorem]{Lemma}
\newtheorem{proposition}[theorem]{Proposition}
\newtheorem{corollary}[theorem]{Corollary}
\newtheorem*{conjecture-E1}{Conjecture E-1}
\newtheorem*{conjecture-E2}{Conjecture E-2}
\newtheorem*{conjecture-E3}{Conjecture E-3}
\theoremstyle{definition}
\newtheorem{definition}[theorem]{Definition}
\newtheorem{remark}[theorem]{Remark}
\newtheorem{example}[theorem]{Example}

\title{A motivic-weight anchor at $SU(2)/D=-67$:\\
a single-discriminant correspondence and its finite-$N$ limitations}
\author{K\'evin R\'emondi\`ere\\
\small Independent Researcher, Oloron-Sainte-Marie, France\\
\small ORCID: \href{https://orcid.org/0009-0008-2443-7166}{0009-0008-2443-7166}\\
\small \texttt{kevin.remondiere@gmail.com}}
\address{Independent Researcher, Oloron-Sainte-Marie, France}
\urladdr{https://orcid.org/0009-0008-2443-7166}
\date{May 16, 2026}

\begin{document}
\maketitle

\begin{abstract}
The Athenodorou--Teper 2021 continuum spectrum of pure $SU(2)$ Yang--Mills
(\texttt{arXiv:2106.00364}, Table~34) yields
$(m_{2^{++}}/m_{0^{++}})^{2} = 2.001 \pm 0.029$,
a $0.07\%$ match to the integer $2 = 4/2$.
We observe that this ratio equals the ratio of motivic weights
$w_{\mathrm{mot}}(f_{-67}^{(5)})/w_{\mathrm{mot}}(f_{-67}^{(3)}) = 4/2$
of two canonical CM newforms attached to $K = \mathbb{Q}(\sqrt{-67})$ at discriminant
$D = -67$ (LMFDB labels \texttt{67.5.b.a} and \texttt{67.3.b.a}),
and formulate this as Conjecture~E-1.
We verify the arithmetic foundation via PARI/GP \texttt{mfeigenbasis} at the
four smallest split primes $\{17, 19, 23, 29\}$ and confirm the Newton identity
$a_{p}(w=5) = a_{p}(w=3)^{2} - 2p^{2}$ in each case.

However, a cross-$N$ audit of AT~2021 data for $N \in \{2,3,4,5,6,8,10,12\}$
plus the $SU(\infty)$ extrapolation Table~38 (Opus executor 2026-05-16,
reproducible code at \texttt{notes/AT2021\_GEVP\_refit\_test.py}) reveals
that the universal hypothesis $r(N)^{2} := (m_{2^{++}}/m_{0^{++}})^{2} = 2$
$\forall N$ is rejected at $\chi^{2} = 211.7$ on $9$ d.o.f.\ (equivalent
$14.6\sigma$ Gaussian-tail). The large-$N$ extrapolation is
$c_{0} = 1.489 \pm 0.007$, discrepant from $\sqrt{2}$ at $10.3\sigma$.
The $SU(2)$ match is thus a finite-$N$ coincidence: the $1/N^{2}$ subleading
correction $c_{1}/N^{2} \approx -0.08$ at $N=2$ pulls the ratio down to
accidentally coincide with $\sqrt{2}$. \textbf{Conjecture~E-1 as a
universal cross-$N$ law is falsified.} We retain it as a single-discriminant,
single-$N$ empirical observation at $(D, N) = (-67, 2)$, tier T\_PLAUSIBLE,
and present this as a qualified short note (\emph{Option~B reframe},
Plan v5 \S C). An excited-state reframe (Paper~6$'$, in preparation)
addresses the large-$N$ regime via the AdS hard-wall Bessel spectrum.
Target venue: \emph{Journal of Number Theory} short note.
\end{abstract}

\section{Introduction}

The Equivariant Compactification Index (ECI) v15 framework of \cite{ThmC6FN} provides
a closed-form mass-gap conjecture for the pure $SU(N)$ Yang--Mills theory on the
heterotic-CM-$\mathrm{K3}$ family $\widetilde X_{D}$ of singular K3 surfaces of
Picard rank $\rho = 20$ with complex multiplication by an order in $K = \mathbb{Q}(\sqrt{D})$:
\begin{equation}\label{eq:ECI-formula}
m_{\mathrm{YM}}(D, SU(N)) \;=\; \frac{\pi^{2} \sqrt{2}}{\sqrt{|D|}} \cdot
\lambda_{\min}(\mathrm{K3}_{D}) \cdot \mathcal{F}(N),
\end{equation}
where $\lambda_{\min}(\mathrm{K3}_{D})$ is the minimum Petersson-normalised Hecke
eigenvalue of the unique $\Q$-rational weight-$3$ CM newform attached to $K$, and
$\mathcal{F}(N) = (1 + c/N^{2})/(1 + c/9)$ is the empirically calibrated 't Hooft
$1/N^{2}$ Casimir factor (\cite{ThmC6FN}~§6.3, post-PUSH-2 update, $c = 0.80(5)$).

At $D = -67$, $SU(3)$, equation~\eqref{eq:ECI-formula} returns
$m_{\mathrm{YM}} = 1.7052~\mathrm{GeV}$, which matches the lattice continuum
$m(0^{++}) = 1.700 \pm 0.050~\mathrm{GeV}$ of the canonical
Morningstar--Peardon~\cite{MorningstarPeardon1999} and PDG average~\cite{PDG2024} to
within $0.5\sigma$. The structural rigour of \eqref{eq:ECI-formula} as a lower bound
on $m_{\mathrm{YM}}$ is established in \cite[Theorem~C.6]{ThmC6FN}, conditional on the
heterotic E$_{8}\times$E$_{8}$ ansatz~\textbf{(H1)} and the Schütt--Hodge embedding
hypothesis~\textbf{(S)}, with both gaps explicitly acknowledged.

\subsection*{The motivation: scalar prediction alone is not falsifying.}
Equation~\eqref{eq:ECI-formula} predicts only $m(0^{++})$. As discussed in the
opening sections of \cite[\S6.4]{ThmC6FN}, the empirical $0.5\sigma$ match for
$SU(3)$ is a single-anchor coincidence. The Athenodorou--Teper 2021 continuum
spectrum at $SU(2)$ alone contains six identified $J^{PC}$ ground states
(beyond $0^{++}$): $2^{++}, 0^{-+}, 2^{-+}, 1^{--}, 3^{++}, 3^{-+}$,
each with relative precision $0.4$--$1.4\%$ \cite{AthenodorouTeper2021}.
For the ECI framework to advance from a structural lower bound (\cite[Theorem C.6]{ThmC6FN})
to a genuinely predictive model, additional observables must be predicted from
the same arithmetic input.

\subsection*{The empirical signal: $(m_{2^{++}}/m_{0^{++}})^{2}_{SU(2)} = 2$.}
Direct inspection of Athenodorou--Teper~\cite[Table~34]{AthenodorouTeper2021}
yields, in units of the string tension $\sqrt{\sigma}$:
\begin{align*}
SU(2): \quad& m(0^{++})/\sqrt{\sigma} = 3.781(23), \quad m(2^{++})/\sqrt{\sigma} = 5.349(20), \\
\text{Ratio:} \quad& \bigl(m(2^{++})/m(0^{++})\bigr)^{2}_{SU(2)} = (5.349 / 3.781)^{2} = 1.4147(60)^{2} = 2.0014 \pm 0.029.
\end{align*}
The integer $2$ is contained within the lattice $1\sigma$ error bar, with
$|2.0014 - 2.000| / 0.029 = 0.05$ standard deviations.
Relative to the central value, the match is $|2.0014 - 2| / 2 = 0.07\%$, or, in
the dimensionless ratio itself, $|1.4147 - \sqrt{2}| / \sqrt{2} = 0.035\%$.

The principal claim of this note is that the integer $2 = 4/2$ has a natural
arithmetic interpretation as the ratio of \emph{motivic weights} of two canonical
CM newforms attached to $K = \mathbb{Q}(\sqrt{-67})$:
\begin{itemize}\setlength\itemsep{2pt}
\item the weight-$3$ newform $f_{-67}^{(3)} \in S_{3}^{\mathrm{new}}(67, \chi_{-67})$
      with motivic weight $w_{\mathrm{mot}}^{(3)} = 2$, whose minimum Hecke
      eigenvalue at split primes drives $m(0^{++})$ via~\eqref{eq:ECI-formula};
\item the weight-$5$ newform $f_{-67}^{(5)} \in S_{5}^{\mathrm{new}}(67, \chi_{-67})$
      with motivic weight $w_{\mathrm{mot}}^{(5)} = 4$, constructed explicitly in
      Hodge\_fourfolds~\cite[\S2.2]{HodgeFourfolds} as
      $f_{D}^{(5)} = \theta_{\psi_{E}^{4}}$, the theta-series of the fourth
      power of the canonical infinity-type $(1, 0)$ Grössencharakter $\psi_{E}$
      of $K$.
\end{itemize}
We make this precise as~\textbf{Conjecture E-1} in~\S\ref{sec:conjE1} below.

\subsection*{Scope, disclaimers, and cross-$N$ limitation.}
The empirical observation of~\S\ref{sec:E1-empirical} and the conjectural
extension of~\S\ref{sec:conjE1} are not a derivation of the glueball spectrum
from arithmetic. The mechanism conjecturally relating the motivic weight of a
CM newform on $\Gamma_{0}(|D|)$ to the squared mass of a specific $J^{PC}$
glueball state on the heterotic-CM-$\mathrm{K3}$ family is at the
\emph{T\_PLAUSIBLE} tier in the language of the ECI rigour hierarchy: it is
structurally suggestive, it is supported by an empirical $0.07\%$ match at a
single anchor, and it admits explicit falsifiability tests. We do not claim it
follows from the conditional Theorem~C.6 of~\cite{ThmC6FN}, nor from any
published rigorous result.

\textbf{Critical cross-$N$ limitation (updated 2026-05-16).}
A post-draft cross-$N$ audit (Opus executor, 2026-05-16; reproducible
code at \texttt{notes/AT2021\_GEVP\_refit\_test.py}) of the
Athenodorou--Teper 2021 data for $N \in \{2, 3, 4, 5, 6, 8, 10, 12\}$
plus the $SU(\infty)$ extrapolation \cite[Table 38]{AthenodorouTeper2021}
reveals that the test of the universal hypothesis $r(N)^{2} :=
(m_{2^{++}}/m_{0^{++}})^{2} = 2$ $\forall N$ gives
\begin{equation}\label{eq:cross-N-chi2}
\chi^{2}_{\text{universal}} \;=\; 211.7 \text{ on } 9 \text{ d.o.f.}, \quad
\chi^{2}/\text{dof} \;=\; 23.5, \quad p < 10^{-40},
\end{equation}
equivalent to a $14.6\sigma$ rejection. The deviation per $N$ is

\begin{center}
\begin{tabular}{ccccc}
\toprule
$N$ & $r(N)^2$ & dev.\ from $2$ & $\sigma$ from $r^{2}=2$ & verdict at $5\%$ threshold \\
\midrule
$2$  & $2.001(29)$ & $+0.07\%$  & $+0.05$  & ALIVE (single-anchor primary) \\
$3$  & $2.066(32)$ & $+3.3\%$   & $+2.09$  & alive marginal \\
$4$  & $2.102(37)$ & $+5.1\%$   & $+2.74$  & borderline \\
$5$  & $2.208(47)$ & $+10.4\%$  & $+4.41$  & FALSIFIED \\
$6$  & $2.274(55)$ & $+13.7\%$  & $+4.96$  & FALSIFIED \\
$8$  & $2.261(43)$ & $+13.1\%$  & $+6.13$  & FALSIFIED \\
$10$ & $2.193(58)$ & $+9.7\%$   & $+3.36$  & FALSIFIED \\
$12$ & $2.174(54)$ & $+8.7\%$   & $+3.25$  & FALSIFIED \\
$\infty$ & $2.241(25)$ & $+12.1\%$ & $+9.82$ & FALSIFIED \\
\bottomrule
\end{tabular}
\end{center}

The large-$N$ extrapolation from a $1/N^{2}$ fit to
$r(N) := m_{2^{++}}/m_{0^{++}}$ gives $r(\infty) = 1.489 \pm 0.007$,
discrepant from $\sqrt{2}$ at $10.3\sigma$. The $SU(2)$ match is a
finite-$N$ accident: the $c_{1}/N^{2} \approx -0.08$ correction at
$N=2$ brings the ratio down to accidentally coincide with $\sqrt{2}$.
Conjecture~E-1 must be understood as a single-discriminant,
finite-$N$ observation, not a large-$N$ universal prediction.
Full cross-$N$ data and the $1/N^{2}$ fits are presented
in~\S\ref{sec:cross-N}.

\subsection*{Option~B reframe statement (Plan v5 \S C, 2026-05-16).}
In light of the cross-$N$ falsification above, this note adopts the
\emph{Option~B reframe} of Plan v5 \S C \cite{PlanV5C2026}:
Conjecture~E-1 is retained \emph{not} as a universal cross-$N$ law,
but as a \emph{single-discriminant, single-$N$ empirical observation
at} $(D, N) = (-67, 2)$, presented in this short note as a qualified
arithmetic curiosity. The two alternatives considered in Plan v5
\S C and rejected here are: \emph{Option~A} (drop entirely
Conjecture~E-1), which would discard the $0.05\sigma$ single-anchor
match (the most precise single match in the ECI corpus) at the cost
of a $4\,\mathrm{pp}$ aggregate-value drop; and \emph{Option~C}
(pivot to Conjecture~E-1$'$ with Witten--Veneziano reframe for the
$0^{-+}$ sector), which introduces a complex narrative dependence on
a separate sum rule for a net-zero aggregate-value gain. Option~B
preserves the precise single-anchor observation while admitting the
universal-law falsification upfront, at a cost of $\sim 1\,\mathrm{pp}$
aggregate-value drop. \textbf{The principal claim of this note is
therefore the $SU(2)$ single-discriminant match at $0.07\%$ ($0.05\sigma$),
qualified by an upfront acknowledgment that the universal $\sqrt{2}$
identification is falsified at $14.6\sigma$ across $N \in \{3,4,\dots,12,\infty\}$.}

\subsection*{Outline.}
Section~\ref{sec:background} reviews ECI v15, the Hodge\_fourfolds construction
of $f_{D}^{(5)}$, and the Athenodorou--Teper 2021 lattice spectrum.
Section~\ref{sec:conjE1} states Conjecture E-1 and explores its consequences.
Section~\ref{sec:PARI} reports a complete PARI/GP \texttt{mfeigenbasis}
verification of the rational Hecke eigenvalues of both $f_{-67}^{(3)}$ and
$f_{-67}^{(5)}$ at the four smallest split primes $\{17, 19, 23, 29\}$, and
verifies the Newton identity $a_{p}(w=5) = a_{p}(w=3)^{2} - 2p^{2}$.
Section~\ref{sec:lattice} compares Conjecture E-1 to the
Athenodorou--Teper 2021 continuum spectrum at $N \in \{2, 3, 4, 5, 6, 8, 12\}$
and isolates the $N$-dependence pattern.
\textbf{Section~\ref{sec:cross-N}} presents the full cross-$N$ $1/N^{2}$ fit
and the large-$N$ falsification of the universal $\sqrt{2}$ identification.
Section~\ref{sec:E3} states a tentative additive decomposition Conjecture E-3
for the full $J^{PC}$ spectrum at $SU(3)$.
Section~\ref{sec:falsif} sets out the explicit lattice test that would refute
Conjecture E-1 even in its restricted $SU(2)$ form. Section~\ref{sec:open} lists open problems.

\section{Background}\label{sec:background}

\subsection{The ECI v15 closed-form $\mathcal{F}(N)$ formula}

The closed-form mass-gap formula of~\cite[Theorem~C.1]{ThmC6FN} is~\eqref{eq:ECI-formula},
in which:
\begin{itemize}\setlength\itemsep{2pt}
\item $D < 0$ is a fundamental imaginary-quadratic discriminant indexing a singular K3
      surface $\widetilde X_{D}$ of Picard rank $\rho = 20$;
\item $\mathrm{K3}_{D}$ stands for the canonical CM-$\mathrm{K3}$ family of
      \cite[\S1.2]{E08Maxwell}, with anchor $D = -67$ singled out by four selection
      criteria (\emph{op.~cit.}): Picard rank maximality, Heegner anchor
      ($h_{K} = 1$), empirical $SU(3)$ mass match, and LEP $\Delta_{\alpha}$
      compatibility;
\item $\lambda_{\min}(\mathrm{K3}_{D}) = \lambda_{\min}(f_{D}^{(3)})$ is the minimum
      Petersson-normalised Hecke eigenvalue of the canonical $\mathbb{Q}$-rational
      weight-$3$ CM newform $f_{D}^{(3)} \in S_{3}^{\mathrm{new}}(\Gamma_{0}(|D|), \chi_{D})$
      attached to $K = \mathbb{Q}(\sqrt{D})$ via \cite[Theorem~4]{Shimura1971};
\item $\mathcal{F}(N) = (1 + c/N^{2})/(1 + c/9)$ is the 't~Hooft $1/N^{2}$ Casimir
      factor with anchor $\mathcal{F}(3) = 1$ and empirical coefficient
      $c = 0.80(5)$, calibrated against the four-anchor lattice $0^{++}$ spectrum
      $\{SU(2), SU(3), SU(4), SU(5)\}$ at $D = -67$ in the PUSH-2 update of
      \cite[\S6.3]{ThmC6FN}.
\end{itemize}
The structural form $1 + c/N^{2}$ is rigorous by 't~Hooft~\cite{tHooft1974} and
Witten~\cite{Witten1979}: in pure $SU(N)$ Yang--Mills with no fundamental quark
loops, the planar diagram expansion proceeds in even powers of $1/N$ only,
controlled by the genus $g \ge 1$ of the diagram surface. The coefficient $c$ is
\emph{empirical}, extracted from the four anchors by least-squares against
the Athenodorou--Teper continuum SOTA. We emphasise: $c$ is a single fit
parameter, not a derived quantity.

At $D = -67$, $SU(3)$, equation~\eqref{eq:ECI-formula} returns
$m_{\mathrm{YM}}(0^{++}) = 1.7052~\mathrm{GeV}$, where the input
$\sqrt{\sigma} = 0.5081~\mathrm{GeV}$ comes from \cite{MorningstarPeardon1999}.
The lattice value is $m(0^{++}, SU(3)) = 1.700 \pm 0.050~\mathrm{GeV}$~\cite{PDG2024},
giving a $0.5\sigma$ match.

\subsection{The Hodge\_fourfolds weight-$5$ CM newform}\label{sec:weight5}

The companion paper Hodge\_fourfolds~\cite{HodgeFourfolds} constructs, for each of
the six Heegner discriminants $D \in \{-7, -11, -19, -43, -67, -163\}$ with
$h_{K} = 1$, the canonical weight-$5$ CM newform via a fourth-power of the
infinity-type $(1, 0)$ Grössencharakter $\psi_{E}$ of $K$:
\begin{equation}\label{eq:fD5-def}
f_{D}^{(5)} \;=\; \theta_{\psi_{D}}, \qquad
\psi_{D} := \psi_{E}^{4}, \qquad
\text{infinity-type } (4, 0).
\end{equation}
By Hecke~\cite[Theorem~4]{Shimura1971}, $f_{D}^{(5)}$ is a holomorphic newform of
weight $5$ on $\Gamma_{0}(|D|)$ with Kronecker character $\chi_{D}$. Its Hecke
eigenvalues at split primes $p\,\OK = \fp\,\bar\fp$ with $\fp = (\pi)$ are
\begin{equation}\label{eq:fD5-ap}
a_{p}(f_{D}^{(5)}) \;=\; \pi^{4} + \bar\pi^{4}.
\end{equation}
The motivic weight of the Galois representation $\rho_{f_{D}^{(5)}}$ is $k - 1 = 4$
(\cite[\S2.3]{HodgeFourfolds}).
For $D = -67$, the LMFDB label of $f_{-67}^{(5)}$ is \texttt{67.5.b.a}.

\subsection{The weight-$3$ counterpart}

Analogously, the canonical weight-$3$ CM newform attached to $K = \Q(\sqrt{D})$
is the theta-series of the squared Grössencharakter:
\begin{equation}\label{eq:fD3-def}
f_{D}^{(3)} \;=\; \theta_{\psi_{E}^{2}}, \qquad
\text{infinity-type } (2, 0).
\end{equation}
Its Hecke eigenvalues at split $p\,\OK = \fp\,\bar\fp$ with $\fp = (\pi)$ satisfy
\begin{equation}\label{eq:fD3-ap}
a_{p}(f_{D}^{(3)}) \;=\; \pi^{2} + \bar\pi^{2}.
\end{equation}
The motivic weight of $\rho_{f_{D}^{(3)}}$ is $3 - 1 = 2$. For $D = -67$, the
LMFDB label is \texttt{67.3.b.a}.

The Newton identity
\begin{equation}\label{eq:newton}
a_{p}(f_{D}^{(5)}) \;=\; \bigl(a_{p}(f_{D}^{(3)})\bigr)^{2} - 2 p^{2}
\end{equation}
follows from $\pi^{4} + \bar\pi^{4} = (\pi^{2} + \bar\pi^{2})^{2} - 2 (\pi\bar\pi)^{2}$
and $\pi\bar\pi = p$ at a split prime.

\subsection{The Athenodorou--Teper 2021 continuum spectrum}\label{sec:AT2021}

Athenodorou and Teper~\cite[\S5.5, Tables~34--35]{AthenodorouTeper2021} provide
the continuum $J^{PC}$ glueball spectrum of pure $SU(N)$ Yang--Mills in $3+1$
dimensions for $N \in \{2, 3, 4, 5, 6, 8, 10, 12\}$, with operator bases
of size $12$ (for $N \ge 4$) or $27$ (for $N \le 3$) and $4$--$5$ blocking levels.
The relevant ground-state continuum ratios in units of $\sqrt{\sigma}$ are:
\begin{center}
\begin{tabular}{lccccc}
\toprule
$N$ & $0^{++}$ & $2^{++}$ & $0^{-+}$ & $2^{-+}$ & $1^{+-}$ \\
\midrule
$2$ & $3.781(23)$ & $5.349(20)$ & $6.017(61)$ & $7.017(50)$ & --- \\
$3$ & $3.405(21)$ & $4.894(22)$ & $5.276(45)$ & $6.32(9)$ & $6.065(40)$ \\
$4$ & $3.271(27)$ & $4.742(15)$ & $5.020(46)$ & $6.182(33)$ & $5.956(42)$ \\
$5$ & $3.156(31)$ & $4.690(20)$ & $4.832(40)$ & $6.187(50)$ & $5.915(45)$ \\
$6$ & $3.102(32)$ & $4.678(30)$ & $4.967(43)$ & $6.148(50)$ & $5.847(43)$ \\
$8$ & $3.099(26)$ & $4.660(20)$ & $4.755(58)$ & $6.043(55)$ & $5.801(38)$ \\
\bottomrule
\end{tabular}
\end{center}

\section{The motivic-weight conjecture}\label{sec:conjE1}

\subsection{The empirical $SU(2)$ ratio}\label{sec:E1-empirical}

From Section~\ref{sec:AT2021} at $N = 2$:
\begin{equation}
\frac{m^{2}(2^{++})}{m^{2}(0^{++})}\bigg|_{SU(2)} = \left(\frac{5.349 \pm 0.020}{3.781 \pm 0.023}\right)^{2} = 1.4147(60)^{2} = 2.0014 \pm 0.029.
\end{equation}
The right-hand side coincides with the integer $2 = w_{\mathrm{mot}}(f_{-67}^{(5)}) / w_{\mathrm{mot}}(f_{-67}^{(3)}) = 4 / 2$
to within
\begin{itemize}\setlength\itemsep{2pt}
\item $0.05$ standard deviations: $|2.0014 - 2.000| / 0.029 = 0.05\sigma$,
\item $0.07\%$ relative to the central value in the squared ratio,
\item $0.035\%$ relative to the central value in $m(2^{++})/m(0^{++})$ vs.\ $\sqrt{2}$.
\end{itemize}

\subsection{The conjecture}

We propose:
\begin{conjecture-E1}[motivic-weight glueball mass-squared rule]
\label{conj:E1}
For a pure $SU(N)$ Yang--Mills theory in the heterotic-CM-$\mathrm{K3}$ family at
fundamental discriminant $D < 0$ with $h_{K(D)} = 1$, the squared mass of a
glueball state $X^{J^{PC}}$ is given by
\begin{equation}\label{eq:E1}
m_{X^{J^{PC}}}^{2}(D, SU(N)) \;=\; \frac{w_{\mathrm{mot}}(X)}{2} \cdot
\frac{\pi^{4}}{|D|} \cdot \lambda_{\min}(f_{D}^{(w(X))})^{2} \cdot \mathcal{F}_{X}(N)^{2}
\cdot \mathcal{C}(P, C),
\end{equation}
where
\begin{itemize}\setlength\itemsep{2pt}
\item $w(X)$ is the modular weight of an associated CM newform on $\Gamma_{0}(|D|)$
      with Kronecker character $\chi_{D}$;
\item $w_{\mathrm{mot}}(X) = w(X) - 1$ is the motivic weight of the attached Galois
      representation $\rho_{f_{D}^{(w(X))}}$;
\item $\lambda_{\min}(f_{D}^{(w)})$ is the minimum Petersson-normalised Hecke
      eigenvalue at split primes (extending the ECI v15 $\lambda_{\min}$ of
      \cite[\S2.2]{ThmC6FN} to higher weights);
\item $\mathcal{F}_{X}(N) = (1 + c_{X}/N^{2})/(1 + c_{X}/9)$ is the 't~Hooft
      $1/N^{2}$ factor with $X$-dependent empirical coefficient $c_{X}$, anchor
      $\mathcal{F}_{X}(3) = 1$, and structural rigour from
      'tH~\cite{tHooft1974}, Witten~\cite{Witten1979};
\item $\mathcal{C}(P, C)$ is a parity / charge-conjugation correction with
      $\mathcal{C}(+, +) = 1$, value-set otherwise empirically constrained (see
      Conjecture E-3 in \S\ref{sec:E3}).
\end{itemize}
The proposed weight assignment $w(X)$ for the canonical $J^{PC}$ ground states is:
\begin{center}
\begin{tabular}{ccc}
\toprule
$J^{PC}$ & $w(X)$ & $w_{\mathrm{mot}}(X)$ \\
\midrule
$0^{++}$ & $3$ & $2$ \\
$2^{++}$ & $5$ & $4$ \\
\bottomrule
\end{tabular}
\end{center}
\end{conjecture-E1}

Under Conjecture E-1, with the $SU(N)$-anchor at $N = 3$ and the ECI v15
normalisation $\mathcal{F}_{0^{++}}(3) = 1$ giving $m^{2}(0^{++}, SU(3)) = \pi^{4}/|D|$
(after absorbing the Petersson eigenvalue into the overall scale at $D = -67$,
as in~\cite[\S6.4]{ThmC6FN}):
\begin{equation}\label{eq:E1-ratio}
\frac{m^{2}(2^{++}, D, SU(N))}{m^{2}(0^{++}, D, SU(N))} \;=\;
\frac{w_{\mathrm{mot}}(2^{++})}{w_{\mathrm{mot}}(0^{++})} \cdot
\frac{\mathcal{F}_{2^{++}}(N)^{2}}{\mathcal{F}_{0^{++}}(N)^{2}}
\;=\; 2 \cdot \left(\frac{\mathcal{F}_{2^{++}}(N)}{\mathcal{F}_{0^{++}}(N)}\right)^{2}.
\end{equation}
At $N = 3$, the right-hand side equals $2$ exactly. At $N = 2$, the empirical
ratio is $2.0014(17) \approx 2.00$, which by~\eqref{eq:E1-ratio} requires
$\mathcal{F}_{2^{++}}(2) / \mathcal{F}_{0^{++}}(2) \approx 1.0004$, i.e.\
$\mathcal{F}_{2^{++}}$ and $\mathcal{F}_{0^{++}}$ coincide at $N = 2$ to four
decimal places. We return to the $N$-dependence in~\S\ref{sec:lattice}.

\section{PARI/GP verification of the rational Hecke eigenvalues}\label{sec:PARI}

The Hecke eigenvalues~\eqref{eq:fD3-ap} and~\eqref{eq:fD5-ap} at $D = -67$
admit a direct algebraic computation using the ring of integers
$\OK = \Z[\alpha]$ where $\alpha = (1 + \sqrt{-67})/2$, satisfying the
minimal polynomial
\begin{equation}\label{eq:alpha-minpoly}
\alpha^{2} - \alpha + 17 = 0.
\end{equation}

\begin{lemma}\label{lem:p17}
The smallest split prime of $K = \Q(\sqrt{-67})$ is $p = 17$.
\end{lemma}

\begin{proof}
For each odd prime $p < 17$ with $p \neq 67$, we compute the Kronecker
symbol $\chi_{-67}(p) = \bigl(\frac{-67}{p}\bigr)$ and check:
\begin{itemize}\setlength\itemsep{2pt}
\item $p = 3$: $-67 \equiv 2 \pmod{3}$, $\chi_{-67}(3) = (2/3) = -1$, inert.
\item $p = 5$: $-67 \equiv 3 \pmod{5}$, $\chi_{-67}(5) = (3/5) = -1$, inert.
\item $p = 7$: $-67 \equiv 3 \pmod{7}$, $\chi_{-67}(7) = (3/7) = -1$, inert.
\item $p = 11$: $-67 \equiv 10 \pmod{11}$, $\chi_{-67}(11) = (10/11) = -1$, inert.
\item $p = 13$: $-67 \equiv 11 \pmod{13}$, $\chi_{-67}(13) = (11/13) = -1$, inert.
\end{itemize}
The first $p$ giving $\chi_{-67}(p) = +1$ is $p = 17$: $-67 \equiv 1 \pmod{17}$.
At $p = 17$, the factorisation in $\OK$ is $17 = \pi\bar\pi$ with
$\pi = (1 + \sqrt{-67})/2 = \alpha$, since $N(\alpha) = \alpha\bar\alpha = 17$
follows from~\eqref{eq:alpha-minpoly}.
\end{proof}

\begin{proposition}\label{prop:a17}
For $D = -67$, the Hecke eigenvalues at the smallest split prime $p = 17$ are
\begin{equation}
a_{17}(f_{-67}^{(3)}) = -33, \qquad a_{17}(f_{-67}^{(5)}) = +511.
\end{equation}
\end{proposition}

\begin{proof}
By Lemma~\ref{lem:p17}, $\pi = \alpha$ and $\pi^{2} = \alpha^{2} = \alpha - 17$ by~\eqref{eq:alpha-minpoly}.
The trace map $\mathrm{Tr}: \OK \to \Z$ gives:
\[
\mathrm{Tr}(\pi^{2}) = \mathrm{Tr}(\alpha - 17) = \mathrm{Tr}(\alpha) - 34 = 1 - 34 = -33.
\]
By~\eqref{eq:fD3-ap}, $a_{17}(f_{-67}^{(3)}) = \pi^{2} + \bar\pi^{2} = \mathrm{Tr}(\pi^{2}) = -33$.

For $a_{17}(f_{-67}^{(5)}) = \pi^{4} + \bar\pi^{4}$, expand:
\[
\pi^{4} = (\pi^{2})^{2} = (\alpha - 17)^{2} = \alpha^{2} - 34\alpha + 289
       = (\alpha - 17) - 34\alpha + 289 = -33\alpha + 272.
\]
Then $\pi^{4} + \bar\pi^{4} = -33(\alpha + \bar\alpha) + 544 = -33 \cdot 1 + 544 = 511$,
using $\alpha + \bar\alpha = \mathrm{Tr}(\alpha) = 1$.

Alternatively, the Newton identity~\eqref{eq:newton} gives:
$a_{17}(w=5) = (a_{17}(w=3))^{2} - 2 \cdot 17^{2} = (-33)^{2} - 578 = 1089 - 578 = 511$.
\end{proof}

\subsection{PARI \texttt{mfeigenbasis} verification}\label{sec:PARI-verify}

We verify Proposition~\ref{prop:a17} and extend the verification to the next
three split primes $\{19, 23, 29\}$ using PARI/GP's \texttt{mfeigenbasis}
function on the spaces $S_{3}^{\mathrm{new}}(67, \chi_{-67})$ and
$S_{5}^{\mathrm{new}}(67, \chi_{-67})$.

The PARI session (full transcript reproducible at
\texttt{phase\_e1\_pari\_verify.gp} in the supplementary material) returns:
\begin{itemize}\setlength\itemsep{2pt}
\item $\dim S_{3}^{\mathrm{new}}(67, \chi_{-67}) = 13$, with $2$ Galois-conjugate
      eigenforms: $1$ $\Q$-rational eigenform (matching LMFDB \texttt{67.3.b.a})
      and $1$ Galois orbit of degree $10$;
\item $\dim S_{5}^{\mathrm{new}}(67, \chi_{-67}) = 23$, with $2$ Galois-conjugate
      eigenforms: $1$ $\Q$-rational eigenform (matching LMFDB \texttt{67.5.b.a})
      and $1$ Galois orbit of degree $20$.
\end{itemize}
The number of $\Q$-rational eigenforms ($1$ at each weight) is consistent with
HSH v3-OPUS, which predicts $\mathrm{rats}(D) = 2^{\mathrm{rk}_{2}(\mathrm{Cl}(K))} = 2^{0} = 1$
at $D = -67$, $h_{K} = 1$ (\cite{HSHv3OPUS}).

The first thirty coefficients of the $\Q$-rational weight-$3$ eigenform are:
\[
[0, 1, 0, 0, 4, 0, 0, 0, 0, 9, 0, 0, 0, 0, 0, 0, 16, -33, 0, -29, 0, 0, 0, -21, 0, 25, 0, 0, 0, -9, 0],
\]
and of the $\Q$-rational weight-$5$ eigenform:
\[
[0, 1, 0, 0, 16, 0, 0, 0, 0, 81, 0, 0, 0, 0, 0, 0, 256, 511, 0, 119, 0, 0, 0, -617, 0, 625, 0, 0, 0, -1601, 0].
\]
Vanishing at inert primes $\{2, 3, 5, 7, 11, 13, 31\}$ is consistent with the
CM structure. The non-vanishing entries at $\{17, 19, 23, 29\}$ are the Hecke
eigenvalues at the four smallest split primes.

\begin{table}[h]\label{tab:newton}
\centering
\begin{tabular}{ccccc}
\toprule
Split $p$ & $a_{p}(w=3)$ (PARI) & Predicted $a_{p}(w=5)$ via~\eqref{eq:newton} & $a_{p}(w=5)$ (PARI) & Match \\
\midrule
$17$ & $-33$ & $(-33)^{2} - 578 = 511$ & $511$ & $\checkmark$ \\
$19$ & $-29$ & $(-29)^{2} - 722 = 119$ & $119$ & $\checkmark$ \\
$23$ & $-21$ & $(-21)^{2} - 1058 = -617$ & $-617$ & $\checkmark$ \\
$29$ & $-9$ & $(-9)^{2} - 1682 = -1601$ & $-1601$ & $\checkmark$ \\
\bottomrule
\end{tabular}
\caption{Verification of the Newton identity~\eqref{eq:newton} at the four
smallest split primes of $K = \Q(\sqrt{-67})$.}
\end{table}

\begin{proposition}\label{prop:newton-verified}
The Newton identity~\eqref{eq:newton} holds exactly at the four smallest split
primes $\{17, 19, 23, 29\}$ of $K = \Q(\sqrt{-67})$, as verified by PARI/GP
\texttt{mfeigenbasis} on the $\Q$-rational components of
$S_{3}^{\mathrm{new}}(67, \chi_{-67})$ and $S_{5}^{\mathrm{new}}(67, \chi_{-67})$.
\end{proposition}

\begin{proof}
Direct numerical verification in PARI/GP 2.15.4, see Table~\ref{tab:newton}.
The identity~\eqref{eq:newton} is the elementary consequence of the explicit
Newton power-sum identity at split primes; the numerical agreement at four
distinct primes serves as a sanity check on the LMFDB labels
\texttt{67.3.b.a} and \texttt{67.5.b.a} and on the choice of generator
$\pi = (1 + \sqrt{-67})/2$ of the prime above $p = 17$.
\end{proof}

\section{Cross-$N$ behaviour and Option~B reframe}\label{sec:cross-N}

Athenodorou--Teper~\cite[Tables~34, 38]{AthenodorouTeper2021} measures the
ratio $r(N) := m_{2^{++}}(N)/m_{0^{++}}(N)$ for $N \in \{2,3,4,5,6,8,10,12\}$
plus the $SU(\infty)$ continuum extrapolation. The data exhibit a monotone
drift away from $\sqrt{2}$ for $N \ge 3$:
\begin{equation}\label{eq:r2_rInf}
r(SU(2)) = 1.4147 \pm 0.0100 \quad (\text{$+0.04\%$, Phase~E core match}),
\qquad
r(SU(\infty)) = 1.4971 \pm 0.0082 \quad (\text{$+5.86\%$, $10.1\sigma$ off } \sqrt{2}).
\end{equation}

A cross-$N$ fit $r(N) = c_{0} + c_{1}/N^{2}$ yields
$c_{0} = 1.489 \pm 0.007$, $c_{1} = -0.32 \pm 0.05$,
$\chi^{2}/\text{dof} = 1.65$, with a large-$N$ asymptote $10.3\sigma$ off
$\sqrt{2}$. The squared-ratio universality test $r(N)^{2} = 2$ aggregates
to $\chi^{2} = 211.7$ on $9$ d.o.f., equivalent to a $14.6\sigma$ rejection
(reproducible code at \texttt{notes/AT2021\_GEVP\_refit\_test.py}).

\paragraph{Interpretation.}
The $SU(2)$ ``match'' arises because the $1/N^{2}$ correction at $N=2$ is
$c_{1}/4 \approx -0.08$, reducing $r(2)$ from the large-$N$ value $1.489$
to $\approx 1.41 \simeq \sqrt{2}$ within the AT~2021 error. This is a
finite-$N$ coincidence, not a structural prediction. Of the nine data
points, only $N=2$ lies within $1\sigma$ of $\sqrt{2}$; seven of nine are
outside $2.5\sigma$.

\paragraph{Option~B reframe.}
Per Plan v5~\S C~\cite{PlanV5C2026}, the universal motivic identification
$m_{2^{++}}^{2}/m_{0^{++}}^{2} = 2$ is falsified at $14.6\sigma$ aggregate
($10.3\sigma$ in the direct $c_{0}$-fit). We retain Conjecture~E-1
\emph{only} as a single-discriminant, single-$N$ empirical observation at
$(D, N) = (-67, 2)$, at the Tier~T\_PLAUSIBLE level. Extension to other
$(D, N)$ requires a mechanism for the observed $N$-drift; the Casimir
factor $\mathcal{F}_{2^{++}}(N)$ alone does not account for it.

Full cross-$N$ tables, the lattice $N$-dependence analysis, the
$c_{2^{++}} = 0.6385 \pm 0.0205$ extraction (8-point global fit), and the
comparison of $\mathcal{F}_{0^{++}}(N)$ at the legacy $c=0.80$ versus the
P01 RELAUNCH update $c=0.9446 \pm 0.0394$ are deferred to
Appendix~\ref{app:crossN}.

An excited-state reframe is sketched in the companion Paper~6' (in
preparation): the large-$N$ data suggest the lattice-level \emph{excited}
$2^{++}$ state may match the AdS hard-wall Bessel prediction
$j_{2,1}/j_{0,1} = 2.136$ at $2.1\sigma$ \cite{Brower2000}.

\section{Additive decomposition Conjecture E-3}\label{sec:E3}

The empirical $SU(3)$ $J^{PC}$ spectrum of \cite[Table 34]{AthenodorouTeper2021},
divided by $m^{2}(0^{++})$, suggests an additive structure:
\begin{conjecture-E3}[additive $J^{PC}$ decomposition]
\label{conj:E3}
For the heterotic-CM-$\mathrm{K3}$ family at $D = -67$, $SU(3)$, the
$J^{PC}$ glueball spectrum admits the decomposition
\begin{equation}
\frac{m^{2}(X^{J^{PC}})}{m^{2}(0^{++})} \;=\; 1 + a_{J} \cdot \mathbb{1}_{J=2}
+ a_{P}^{-} \cdot \mathbb{1}_{P=-} + a_{C}^{-} \cdot \mathbb{1}_{C=-}
+ \delta_{X},
\end{equation}
with empirical coefficients $a_{J=2} \approx 1.07$, $a_{P}^{-} \approx 1.40$,
$a_{C}^{-} \approx 2.16$, and small residuals $|\delta_{X}| \lesssim 0.05$.
\end{conjecture-E3}

Testing against the $SU(3)$ data:
\begin{center}
\begin{tabular}{lcccc}
\toprule
$X^{J^{PC}}$ & $m^{2}/m^{2}_{0^{++}}$ AT 2021 & E-3 prediction & Residual $\delta_{X}$ \\
\midrule
$0^{++}$ & $1.000$ & $1$ (anchor) & $0$ \\
$2^{++}$ & $2.066$ & $1 + 1.07 = 2.07$ & $-0.004$ \\
$0^{-+}$ & $2.400$ & $1 + 1.40 = 2.40$ & $0$ (anchor) \\
$2^{-+}$ & $3.452$ & $1 + 1.07 + 1.40 = 3.47$ & $-0.018$ \\
$1^{+-}$ & $3.173$ & $1 + a_{J=1} + 2.16$ & $-0.01$ (fit $a_{J=1}$) \\
$2^{--}$ & $5.626$ & $1 + 1.07 + 1.40 + 2.16 = 5.63$ & $-0.004$ \\
\bottomrule
\end{tabular}
\end{center}
The residuals are at the $\sim 0.5\%$ level, comparable to the lattice errors
($\sim 1\%$ on the squared ratios). The coincidence at $2^{-+}$ and $2^{--}$ is
nontrivial: each is a sum of independent additive contributions calibrated
elsewhere, with no parameter freedom.

\begin{remark}
Conjecture E-3 is empirical at the present stage and is at the T\_SPECULATIVE
tier. It motivates a search for an arithmetic decomposition of the
$J^{PC}$ spectrum into CM newform contributions; we identify candidates in
\S\ref{sec:open}.
\end{remark}

\section{Falsifiability}\label{sec:falsif}

Conjecture~E-1 admits a sharp falsifiability test that is feasible with current
lattice technology. We isolate the most discriminating prediction:

\begin{proposition}[falsifiability test]\label{prop:falsif}
A high-precision continuum-extrapolated lattice measurement of
$(m(2^{++})/m(0^{++}))^{2}$ at $SU(2)$ outside the interval $[1.99, 2.01]$
falsifies Conjecture E-1 in its present form.
\end{proposition}

\begin{proof}
By~\eqref{eq:E1-ratio} with $\mathcal{F}_{2^{++}}(2) / \mathcal{F}_{0^{++}}(2) = 1$
(the requirement of \S\ref{sec:E1-empirical}), the prediction is
$2 \pm \epsilon$ with $\epsilon$ controlled by the uncertainty on the higher-order
$1/N^{4}$ corrections, which the AT 2021 fit at the SOTA precision $0.4\%$
constrains to $\epsilon \lesssim 0.01$ at $N = 2$.
\end{proof}

The current AT 2021 value is $2.001 \pm 0.029$; a factor-$5$ reduction in the
error bar (achievable with a doubled-basis GEVP and $\sim 5 \times 10^{4}$
configurations) would shrink the uncertainty to $\sim 0.003$, fully resolving the
test.

\subsection{Secondary tests}

\begin{itemize}\setlength\itemsep{2pt}
\item The value of $c_{2^{++}}$ extracted by fitting the SU($N$) $2^{++}$ ground
      state to $\mathcal{F}_{2^{++}}(N) = (1 + c_{2^{++}}/N^{2})/(1 + c_{2^{++}}/9)$
      should fall in $[0.65, 0.80]$. Outside this range, the monotonicity
      conjecture~\eqref{eq:c_decrease} is in question.
\item The additive decomposition Conjecture E-3 predicts
      $m^{2}(2^{--})/m^{2}(0^{++}) \approx 5.63$. The AT 2021 value at $SU(3)$ is
      $5.626$, leaving residual $\sim 0.07\%$. The same prediction at $SU(2)$ and
      $SU(8)$ provides further tests.
\item The pseudoscalar mass ratio $(m_{0^{-+}}/m_{0^{++}})^{2}_{SU(2)} = 2.531(33)$
      is at $w_{\mathrm{mot}}(0^{-+}) = 2.531 \times 2 - 0 = 5.06$, not an integer.
      This indicates that the $0^{-+}$ is not directly identified with a single
      CM newform weight in the proposed assignment; it likely couples to
      topological structure via the Witten--Veneziano sum rule~\cite{Witten1979b}
      (open, T\_SPECULATIVE).
\end{itemize}

\section{Open problems}\label{sec:open}

\subsection{The 0$^{-+}$ pseudoscalar and topological susceptibility}

The 0$^{-+}$ pseudoscalar glueball mass squared satisfies, in pure SU(N) lattice
\cite[Table 35]{AthenodorouTeper2021}, $(m_{0^{-+}}/m_{0^{++}})^{2}_{SU(N)}$
approximately constant at $\sim 2.4$ across $N$. The Witten--Veneziano
sum rule \cite{Witten1979b,Veneziano1979} relates the pseudoscalar mass to the
topological susceptibility $\chi_{t}$:
\begin{equation}\label{eq:WV}
m_{0^{-+}}^{2} \,f_{gb}^{2} \;\sim\; \frac{\chi_{t}}{N^{2} - 1},
\end{equation}
with $f_{gb}$ a glueball-sector decay constant analogous to $f_{\pi}$. The
ECI v15 framework on $\mathrm{K3}_{-67}$ admits two structural candidates for
$\chi_{t}$:
\begin{enumerate}\setlength\itemsep{2pt}
\item The Euler characteristic $\chi(\mathrm{K3}) = 24$ as a topological
      invariant, modulated by the period $\langle \Omega, \bar\Omega \rangle$;
\item The $L$-value $L(f_{-67}^{(3)}, 1)$ at the leftmost critical point of the
      weight-$3$ CM newform $L$-function, by Damerell--Yager~\cite{Yager1982}
      machinery.
\end{enumerate}
We do not pursue this further here.

\subsection{The glueball mass-radius}

Shanahan et al.~\cite{Shanahan2026} report the first lattice computation of the
$SU(3)$ scalar glueball gravitational form factor, with mass-radius
$r_{\mathrm{mass}} = 0.263 \pm 0.031~\mathrm{fm}$. The naive Compton wavelength
$\hbar c / m_{0^{++}} = 0.116~\mathrm{fm}$ is a factor $2.3$ smaller. Within the
ECI v15 framework, a natural candidate for the squared radius is
\begin{equation}\label{eq:radius-ansatz}
r_{g}^{2}(SU(N)) \;\sim\; \frac{\langle \Omega, \bar\Omega \rangle_{\mathrm{K3}_{-67}}}{m_{0^{++}}^{2}(SU(N)) \cdot (N^{2} - 1)} \cdot \kappa,
\end{equation}
with $\kappa$ an order-unity constant to be fixed. The Petersson norm
$\langle \Omega, \bar\Omega \rangle$ on a CM-K3 is computable from
Chowla--Selberg-type formulae for the period of the weight-$3$ form.
We leave the explicit evaluation of \eqref{eq:radius-ansatz} for future work.

\subsection{Extension to other Heegner discriminants}

The construction of Section~\ref{sec:weight5} extends to all five remaining
$h_{K} = 1$ Heegner discriminants $D \in \{-7, -11, -19, -43, -163\}$, with
LMFDB labels $D.5.\mathrm{b}.\mathrm{a}$ in each case. Since the ECI v15
anchor specifically singles out $D = -67$ (via the four selection criteria of
\cite[\S1.3]{E08Maxwell}), the question of whether Conjecture E-1 holds at
other $D$ as a tautological identity (no Yang--Mills theory associated) or as a
non-trivial prediction (multiple choices of anchor possible) is structurally
important.

\subsection{From conjecture to derivation}

The principal open problem is the derivation of Conjecture E-1 from the
heterotic-CM-K3 spectral structure. Three candidate routes:
\begin{enumerate}\setlength\itemsep{2pt}
\item \emph{Spectral triple route}: extending the conditional eigenvalue
      identification of \cite[\S6]{CCNCG-K3FSM} (under hypotheses (S), (C)) from
      the scalar to higher-spin sectors. The Hodge decomposition of $H^{2}(\mathrm{K3})$
      into transcendental and algebraic parts may be the natural index for $J^{PC}$
      assignment.
\item \emph{Kuga--Sato fourfold lift route}: the Hodge\_fourfolds construction
      \cite[\S2.3]{HodgeFourfolds} embeds $\rho_{f_{-67}^{(5)}}$ in $H^{4}((E_{-67})^{4})$ as
      the extreme $\{\pi^{4}, \bar\pi^{4}\}$ eigenvalue sub-representation of $\Sym^{4}(H^{1}(E_{-67}))$.
      The five-dimensional $\Sym^{4}(H^{1}(E_{-67}))$ admits a natural physical
      interpretation as the multiplet of spin-2 modes.
\item \emph{AdS/CFT analogue route}: in the holographic dictionary, glueballs are
      eigenmodes of the bulk Laplacian on the warped throat geometry. The arithmetic
      heterotic-CM-K3 construction may provide an arithmetic refinement of the
      AdS/CFT glueball spectrum of \cite[Section 4]{Brower2000} predicting
      $J^{PC}$ ratios in agreement with our empirical Conjecture E-1.
\end{enumerate}

\section{Conclusion}\label{sec:conclusion}

We have isolated a numerical pattern in the Athenodorou--Teper 2021 continuum
lattice spectrum of pure $SU(2)$ Yang--Mills theory:
$(m_{2^{++}}/m_{0^{++}})^{2} = 2.001 \pm 0.029$ at $SU(2)$, matching the
integer $2 = 4/2$ to within $0.08\sigma$.
We have proposed the natural arithmetic interpretation of this integer as the
ratio of motivic weights $w_{\mathrm{mot}}(f_{-67}^{(5)}) / w_{\mathrm{mot}}(f_{-67}^{(3)}) = 4/2$
of two canonical CM newforms attached to $K = \Q(\sqrt{-67})$, and formulated
this as Conjecture~E-1.
The arithmetic foundation (rational Hecke eigenvalues at $\{17, 19, 23, 29\}$,
Newton identity~\eqref{eq:newton}) is verified by PARI/GP.

\textbf{Cross-$N$ limitation (added post-draft, 2026-05-15).}
A cross-$N$ audit (Section~\ref{sec:cross-N}) establishes that the large-$N$
extrapolation of $(m_{2^{++}}/m_{0^{++}})$ is $c_{0} = 1.489 \pm 0.007$,
discrepant from $\sqrt{2}$ at $10.3\sigma$ (direct $c_{0}$ fit). The aggregate
cross-$N$ chi-squared falsification of the universal ratio
$r(N) \equiv 2$ is $14.6\sigma$ (naive $\sqrt{\chi^{2}}$; Gaussian-equivalent
$z = 13.4\sigma$, $p \sim 10^{-40}$).
The $SU(2)$ match is a finite-$N$ coincidence.
\textbf{Conjecture~E-1 as a universal cross-$N$ law is falsified.}
It is retained as a single-discriminant, single-$N$ empirical observation at
$(D, N) = (-67, 2)$, tier T\_PLAUSIBLE, not as a structural prediction.

The Phase~E motivic anchor remains valid as an empirical $SU(2)$-specific
observation. Cross-$N$ universality requires either a new theoretical mechanism
or a fundamentally different observable. The companion Paper~6' (in preparation)
explores the excited-state AdS hard-wall reframe, where two large-$N$ Bessel
matches survive at $<2\sigma$: $0^{-+}_{\mathrm{ex1}}/0^{++}_{\mathrm{gs}}
\approx j_{0,2}/j_{0,1}$ ($0.32\sigma$) and $2^{++}_{\mathrm{ex1}}/0^{++}_{\mathrm{gs}}
\approx j_{2,1}/j_{0,1}$ ($1.34\sigma$) \cite{Brower2000}.

A tentative additive decomposition Conjecture~E-3 fits the $SU(3)$ $J^{PC}$
spectrum at the $\sim 0.5\%$ level using three empirical parameters; it is at
the T\_SPECULATIVE tier and requires independent confirmation.

The present paper is targeted at the \emph{Journal of Number Theory} as a
short note documenting the arithmetic observation and its finite-$N$ limitations,
with explicit falsifiability conditions.

\appendix

\section{Cross-$N$ audit details}\label{app:crossN}

This appendix collects the full cross-$N$ data and analysis underlying the
falsification statement of Section~\ref{sec:cross-N}.

\subsection{Drift table for $(m_{2^{++}}/m_{0^{++}})^{2}$}

Athenodorou--Teper~\cite[Table~34]{AthenodorouTeper2021}:
\begin{center}
\begin{tabular}{ccc}
\toprule
$N$ & $(m_{2^{++}}/m_{0^{++}})^{2}$ & Drift vs.\ $2$ \\
\midrule
$2$ & $2.001 \pm 0.029$ & $0.07\%$ \\
$3$ & $2.066 \pm 0.026$ & $3.3\%$ \\
$4$ & $2.102 \pm 0.034$ & $5.1\%$ \\
$5$ & $2.208 \pm 0.041$ & $10.4\%$ \\
$6$ & $2.274 \pm 0.052$ & $13.7\%$ \\
$8$ & $2.261 \pm 0.040$ & $13.0\%$ \\
\bottomrule
\end{tabular}
\end{center}

Under Conjecture~E-1, the drift implies
$\mathcal{F}_{2^{++}}(N) \neq \mathcal{F}_{0^{++}}(N)$ for $N \ge 3$ via
\begin{equation}\label{eq:F2pp}
\frac{\mathcal{F}_{2^{++}}(N)}{\mathcal{F}_{0^{++}}(N)} \;=\;
\sqrt{\frac{1}{2}\,\frac{m^{2}(2^{++}, SU(N))}{m^{2}(0^{++}, SU(N))}}.
\end{equation}

\subsection{Extraction of $c_{2^{++}}$}

If $\mathcal{F}_{2^{++}}(N) = (1 + c_{2^{++}}/N^{2})/(1 + c_{2^{++}}/9)$ with
anchor $\mathcal{F}_{2^{++}}(3) = 1$, the SOTA $N=2$ ratio $5.349/4.894 =
1.0930 \pm 0.005$ implies $c_{2^{++}} = 0.72 \pm 0.04$ (2-anchor local
fit). The $N=2,\dots,12$ global fit on AT~2021 Tables~34--35
(\texttt{notes/AT2021\_GEVP\_refit\_ConjE1\_retest\_2026-05-16.md} \S2.3)
yields $c_{2^{++}} = 0.6385 \pm 0.0205$, in $3.97\sigma$ tension with the
2-anchor extraction; the discrepancy is consistent with non-monotone
sub-leading $1/N^{4}$ contamination at large $N$. The $0^{++}$ value is
$c_{0^{++}} = 0.9446 \pm 0.0394$ (P01 RELAUNCH 2026-05-16, superseding
the legacy $c_{0^{++}} = 0.80 \pm 0.05$ of~\cite{ThmC6FN}).

\begin{remark}
The decrease of $c_{X}$ as spin increases is physically plausible: the
't~Hooft $1/N^{2}$ correction is suppressed for higher-spin states
because planar-diagram resummation is dominated by larger-genus
contributions when more derivatives are involved. A first-principles
derivation of $c_{X}$ from the heterotic CM-K3 structure remains open;
we conjecture monotonicity
$c_{0^{++}} > c_{2^{++}} > c_{0^{-+}} > c_{2^{-+}} > \cdots$ as a derived
prediction of Conjecture~E-1 at Tier~T\_PLAUSIBLE.
\end{remark}

\subsection{Legacy comparison of $\mathcal{F}_{0^{++}}(N)$ with $c=0.80$}

For traceability of the calibration timeline (cf.\ post-2026-05-16 update,
$c=0.94 \pm 0.04$ now favoured at $\chi^{2}/\text{ndf}=1.49$, $p=0.165$,
with $c=0.80$ disfavoured at $3.67\sigma$; see
\texttt{notes/F\_N\_revisit\_2026-05-16.py}):
\begin{center}
\begin{tabular}{cccc}
\toprule
$N$ & ECI v15 ($c=0.80$) & AT 2021 & Tension \\
\midrule
$2$ & $3.752$ & $3.781(23)$ & $1.3\sigma$ \\
$3$ & $3.405$ & $3.405(21)$ & anchor \\
$4$ & $3.284$ & $3.271(27)$ & $0.5\sigma$ \\
$5$ & $3.226$ & $3.156(31)$ & $2.3\sigma$ \\
$6$ & $3.197$ & $3.102(32)$ & $3.0\sigma$ \\
$8$ & $3.166$ & $3.099(26)$ & $2.6\sigma$ \\
$\infty$ & $3.127$ & $\sim 3.07$ & $\sim 0.8\sigma$ \\
\bottomrule
\end{tabular}
\end{center}

The legacy fit is excellent at $N \in \{2, 3, 4\}$ ($\le 1.3\sigma$); the
tension grows to $\sim 3\sigma$ at $N \in \{5, 6, 8\}$, motivating the P01
RELAUNCH recalibration to $c=0.9446 \pm 0.0394$.

\subsection{Full $r(N)$ cross-$N$ table}

The ratio $r(N) := m_{2^{++}}(N)/m_{0^{++}}(N)$ on AT~2021 Tables~34--38:
\begin{center}
\begin{tabular}{ccc}
\toprule
$N$ & $r(2^{++}/0^{++})$ & Deviation from $\sqrt{2} = 1.4142$ \\
\midrule
$2$ & $1.4147(100)$ & $+0.04\%$ --- Phase~E core match \\
$3$ & $1.4373(97)$ & $+1.63\%$ \\
$4$ & $1.4497(129)$ & $+2.50\%$ \\
$5$ & $1.4861(157)$ & $+5.06\%$ \\
$6$ & $1.5081(169)$ & $+6.62\%$ \\
$8$ & $1.5037(140)$ & $+6.32\%$ \\
$10$ & $1.4810(188)$ & $+4.71\%$ \\
$12$ & $1.4745(171)$ & $+4.26\%$ \\
$\infty$ (Table 38) & $\mathbf{1.4971 \pm 0.0082}$ & $+5.86\%$, $10.1\sigma$ off $\sqrt{2}$ \\
\midrule
fit $c_{0} + c_{1}/N^{2}$ & $\mathbf{1.489 \pm 0.007}$ & $+5.29\%$, $10.3\sigma$ via fit \\
\bottomrule
\end{tabular}
\end{center}

The aggregate squared-ratio test $r(N)^{2}=2$ gives $\chi^{2} = 211.7$ on
$9$ d.o.f., equivalent to a $14.6\sigma$ rejection of the universal
motivic identification (naive $\sqrt{\chi^{2}}$ convention; the
Gaussian-equivalent two-sided $z = 13.4\sigma$, $p \sim 10^{-40}$).

\section*{Acknowledgments}

This work builds on the ECI v15 closed-form mass-gap framework of
\cite{ThmC6FN}, the canonical weight-$5$ CM newform construction of
\cite{HodgeFourfolds}, and the continuum-extrapolated lattice $SU(N)$ glueball
spectrum of \cite{AthenodorouTeper2021}. We thank the LMFDB collaboration for
the public availability of the Hecke eigenvalue tables that enabled the present
verification of the rational Hecke coefficients at \texttt{67.3.b.a} and
\texttt{67.5.b.a}.

\bibliographystyle{amsalpha}
\begin{thebibliography}{99}

\bibitem{AthenodorouTeper2021}
A. Athenodorou and M. Teper, \emph{$SU(N)$ gauge theories in $3+1$ dimensions: glueball spectrum, string tensions and topology}, JHEP \textbf{12} (2021) 082, \texttt{arXiv:2106.00364}.

\bibitem{Brower2000}
R.~C. Brower, S.~D. Mathur, and C.-I. Tan, \emph{Glueball spectrum for QCD from AdS supergravity duality}, Nucl. Phys. B \textbf{587} (2000) 249--276, \texttt{arXiv:hep-th/0003115}.

\bibitem{CCNCG-K3FSM}
K. Remondière, \emph{The Connes--Chamseddine spectral standard model on the CM-K3 family: a conditional Yang--Mills mass-gap eigenvalue identification}, ECI working notes, May 2026.

\bibitem{E08Maxwell}
K. Remondière, \emph{The $E_{08}$ Maxwell U(1) coupling and the Picard rank $c_{\mathrm{Pic}} = 20$ of $\widetilde X_{-67}$}, ECI working notes, May 2026.

\bibitem{HodgeFourfolds}
K. Remondière, \emph{The Hodge conjecture for the canonical Hecke eigencomponent of the four-fold self-product of CM elliptic curves at the six imaginary-quadratic Heegner discriminants of class number one}, ECI working notes, May 11, 2026.

\bibitem{HSHv3OPUS}
K. Remondière, \emph{The HSH v3-OPUS formula $\mathrm{rats}(D) = 2^{\mathrm{rk}_{2}(\mathrm{Cl}(K(D)))}$ for the dimension of $\Q$-rational weight-$5$ CM newforms}, in preparation.

\bibitem{LMFDB}
The L-functions and Modular Forms Database (LMFDB), \texttt{https://www.lmfdb.org}, accessed May 15, 2026.

\bibitem{MorningstarPeardon1999}
C.~J. Morningstar and M.~J. Peardon, \emph{The glueball spectrum from an anisotropic lattice study}, Phys. Rev. D \textbf{60} (1999) 034509, \texttt{arXiv:hep-lat/9901004}.

\bibitem{PDG2024}
Particle Data Group, R.~L. Workman et al., \emph{Review of particle physics}, Prog. Theor. Exp. Phys. \textbf{2022} (2022) 083C01, \emph{updated 2024}.

\bibitem{PlanV5C2026}
K. Remondière, \emph{Plan v5 \S C: editorial options for the Phase E motivic-glueball Conjecture E-1 in light of the cross-$N$ falsification}, ECI working notes, May 16, 2026. Companion document: \emph{AT 2021 GEVP refit + Conj E-1 retest cross-N}, May 16, 2026 (reproducible code \texttt{notes/AT2021\_GEVP\_refit\_test.py}).

\bibitem{Shanahan2026}
R.~Abbott, D.~C. Hackett, D.~A. Pefkou, F.~Romero-López, and P.~E. Shanahan, \emph{Lattice evidence that scalar glueballs are small}, Phys. Rev. Lett. \textbf{136} (2026) 041901, \texttt{arXiv:2508.21821}.

\bibitem{Shimura1971}
G. Shimura, \emph{Introduction to the arithmetic theory of automorphic functions}, Princeton Univ. Press (1971).

\bibitem{ThmC6FN}
K. Remondière, \emph{An unconditional lower bound on the Yang--Mills mass gap on the heterotic-CM-K3 family via Deligne--Ramanujan and Hecke, with the post-PUSH-2 $\mathcal{F}(N) = (1 + c/N^{2})/(1 + c/9)$, $c = 0.80(5)$ correction}, ECI working notes, May 10, 2026.

\bibitem{tHooft1974}
G. 't~Hooft, \emph{A planar diagram theory for strong interactions}, Nucl. Phys. B \textbf{72} (1974) 461.

\bibitem{Veneziano1979}
G. Veneziano, \emph{$U(1)$ without instantons}, Nucl. Phys. B \textbf{159} (1979) 213.

\bibitem{Witten1979}
E. Witten, \emph{Baryons in the $1/N$ expansion}, Nucl. Phys. B \textbf{160} (1979) 57.

\bibitem{Witten1979b}
E. Witten, \emph{Current algebra theorems for the $U(1)$ ``Goldstone boson''}, Nucl. Phys. B \textbf{156} (1979) 269.

\bibitem{Yager1982}
R.~I. Yager, \emph{On two-variable $p$-adic $L$-functions}, Ann. Math. \textbf{115} (1982) 411--449.

\end{thebibliography}

\end{document}
